\documentclass[10pt]{article}
\usepackage{amsmath,amssymb,amsfonts}
\usepackage{geometry}
\usepackage{xcolor}
\usepackage{fancyhdr}
\usepackage{tcolorbox}
\usepackage{mathpazo}
\usepackage{helvet}

\geometry{
    paperwidth=210mm,
    paperheight=297mm,
    left=58mm,
    right=20mm,
    top=22mm,
    bottom=22mm,
    headheight=14pt,
    headsep=15pt
}

% Hệ màu chuẩn Stewart Calculus
\definecolor{stewartcyan}{RGB}{0, 130, 185}
\definecolor{stewartred}{RGB}{204, 0, 0}

% Hộp đóng khung kết quả viền đỏ chuẩn sách gốc
\newtcolorbox{stewartbox}{
    colback=white,
    colframe=stewartred,
    arc=0mm,
    boxrule=0.8pt,
    left=10pt,
    right=10pt,
    top=7pt,
    bottom=7pt
}

% Tiêu đề đầu trang chuẩn dóng theo cột nội dung chính
\pagestyle{fancy}
\fancyhf{}
\fancyhead[L]{\textsf{\textbf{MỤC 2.7}}\quad \textsf{Đạo hàm và Tốc độ Thay đổi}}
\fancyhead[R]{\textsf{\textbf{145}}}
\renewcommand{\headrulewidth}{0pt}

\setlength{\parindent}{0pt}
\setlength{\parskip}{5pt}

\newcommand{\tombstone}{\hfill{\color{stewartcyan}\rule{1.2ex}{1.2ex}}}

\begin{document}

\noindent
\textbf{(b)}\quad
$\begin{aligned}[t]
f'(a) &= \lim_{h \to 0} \frac{f(a + h) - f(a)}{h} \\[6pt]
&= \lim_{h \to 0} \frac{[(a + h)^2 - 8(a + h) + 9] - [a^2 - 8a + 9]}{h} \\[6pt]
&= \lim_{h \to 0} \frac{a^2 + 2ah + h^2 - 8a - 8h + 9 - a^2 + 8a - 9}{h} \\[6pt]
&= \lim_{h \to 0} \frac{2ah + h^2 - 8h}{h} = \lim_{h \to 0} (2a + h - 8) = 2a - 8
\end{aligned}$

\vspace{8pt}
Để kiểm tra lại kết quả của chúng ta trong phần (a), hãy để ý rằng nếu cho $a = 2$, thì $f'(2) = 2(2) - 8 = -4$.\tombstone

\vspace{12pt}
\noindent
\textbf{\textsf{\color{stewartcyan}VÍ DỤ 5}}\quad Sử dụng Phương trình 5 để tìm đạo hàm của hàm số $f(x) = 1/\sqrt{x}$ tại số $a$ ($a > 0$).

\vspace{4pt}
\noindent
\textbf{\textsf{\color{stewartcyan}LỜI GIẢI}}\quad Từ Phương trình 5 ta có
\[
\begin{aligned}
f'(a) &= \lim_{x \to a} \frac{f(x) - f(a)}{x - a} \\[8pt]
&= \lim_{x \to a} \frac{\dfrac{1}{\sqrt{x}} - \dfrac{1}{\sqrt{a}}}{x - a} = \lim_{x \to a} \frac{\dfrac{1}{\sqrt{x}} - \dfrac{1}{\sqrt{a}}}{x - a} \cdot \frac{\sqrt{x}\sqrt{a}}{\sqrt{x}\sqrt{a}} \\[9pt]
&= \lim_{x \to a} \frac{\sqrt{a} - \sqrt{x}}{\sqrt{ax}(x - a)} = \lim_{x \to a} \frac{\sqrt{a} - \sqrt{x}}{\sqrt{ax}(x - a)} \cdot \frac{\sqrt{a} + \sqrt{x}}{\sqrt{a} + \sqrt{x}} \\[9pt]
&= \lim_{x \to a} \frac{-(x - a)}{\sqrt{ax}(x - a)(\sqrt{a} + \sqrt{x})} = \lim_{x \to a} \frac{-1}{\sqrt{ax}(\sqrt{a} + \sqrt{x})} \\[9pt]
&= \frac{-1}{\sqrt{a^2}(\sqrt{a} + \sqrt{a})} = \frac{-1}{a \cdot 2\sqrt{a}} = -\frac{1}{2a^{3/2}}
\end{aligned}
\]

\vspace{2pt}
Bạn có thể kiểm tra lại rằng việc sử dụng Định nghĩa 4 cũng mang lại kết quả tương tự.\tombstone

\vspace{12pt}
Chúng ta đã định nghĩa tiếp tuyến của đường cong $y = f(x)$ tại điểm $P(a, f(a))$ là đường thẳng đi qua $P$ và có hệ số góc $m$ được cho bởi Phương trình 1 hoặc 2. Vì theo Định nghĩa 4 (và Phương trình 5), giá trị này cũng chính là đạo hàm $f'(a)$, nên giờ đây chúng ta có thể khẳng định như sau:

\vspace{6pt}
\begin{stewartbox}
Đường tiếp tuyến của $y = f(x)$ tại $(a, f(a))$ là đường thẳng đi qua $(a, f(a))$ có hệ số góc bằng $f'(a)$, tức đạo hàm của $f$ tại $a$.
\end{stewartbox}

\vspace{6pt}
Nếu sử dụng dạng điểm -- hệ số góc của phương trình đường thẳng, ta có thể viết phương trình tiếp tuyến của đường cong $y = f(x)$ tại điểm $(a, f(a))$ như sau:
\[
y - f(a) = f'(a)(x - a)
\]

\end{document}